%% parts/p04/universal-reductions.tex --- Phase-4 migration of v1.1 Part V,
%% ch:universal-reductions, source lines 7851--8247.  Body = statements only;
%% all source complete proofs and proof sketches live in the parallel annex.

\chapter[Universal Computation and Reductions]{Universal Computation, Reductions, and Fixed Points}
\label{sec:p04-universal-reductions}

\section{An effective numbering of programs}
\label{sec:p04-effective-numbering}

\begin{hbconvention}{def:p04-program-numbering}{Fixed program numbering}
A numerical-machine description is \emph{normalized} exactly when all of the
following conditions hold:
\begin{enumerate}
\item its tape symbols are numbered
      \(0,\ldots,\abs{\mathcal A_{\mathrm{tape}}}-1\), with
      \(\sqcup\mapsto0\), digit \(0\mapsto1\), and digit \(1\mapsto2\);
\item its states are numbered \(0,\ldots,\abs Q-1\), with
      \(q_0\mapsto0\), \(q_{\mathrm{acc}}\mapsto1\), and
      \(q_{\mathrm{rej}}\mapsto2\), while remaining numbers are assigned
      arbitrarily; and
\item its row list is the lexicographic listing, by numbered key \((q,a)\),
      of the full transition table over its nonhalting domain.
\end{enumerate}
The designated
assignments are part of validity: an ``order of first appearance'' would be
circular because lexicographic row order already presupposes the numbering.

Encode a description tuple \(\vec n\) of length \(k\) as
\[
  \operatorname{pair}_{\N}(k,\operatorname{code}_k(\vec n)).
\]
There is no canonical machine-to-code map, because indices need not be unique.
An invalid code denotes a machine diverging everywhere.  In particular, a
valid numerical-machine code must have
\(\{0,1\}\subseteq\mathcal A_{\mathrm{in}}\), read at designated symbols
\(1,2\), and its three designated states are pairwise distinct.  A source
machine whose start state is halting is represented behaviorally by the
surrogate-start construction in Theorem~\ref{thm:p04-enumeration}.

For a valid code \(e\), \(\varphi_e(x)\) is the numerical output exactly when
the decoded machine accepts with a valid numerical output; it is undefined if
the machine rejects, diverges, or halts accepting with invalid output.  A fixed
wrapper copies a valid result into canonical form and loops after rejection or
invalid acceptance.  The number \(e\) is an \emph{index}.  For
\(\vec x\in\N^k\), \(\varphi_e(\vec x)\) means the behavior of this same
program on \(\operatorname{bin}(\operatorname{code}_k(\vec x))\): it is
arity-dependent application notation, not another numbering or a second unary
coding operation.
\index{program numbering}
\index{index!of a partial computable function}
\end{hbconvention}

\begin{hblemma}{thm:p04-renaming-invariance}{proved=annex:p04-universal.1}{Renaming invariance}
Let \(\mathcal T\) be a numerical-function Turing machine, and let
\(\iota=(\iota_Q,\iota_A)\) bijectively rename its states and tape symbols
while preserving the roles of \(q_0,q_{\mathrm{acc}},q_{\mathrm{rej}}\) and
\(\sqcup,0,1\).  If \(\mathcal T^\iota\) and \(C^\iota\) denote the renamed
machine and configuration, then
\[
  C\longrightarrow_{\mathcal T}C'
  \quad\Leftrightarrow\quad
  C^\iota\longrightarrow_{\mathcal T^\iota}(C')^\iota.
\]
Consequently the two machines, on letterwise corresponding inputs, accept,
reject, or diverge together and have corresponding tapes at every step.
\end{hblemma}

\begin{hbtheorem}{thm:p04-enumeration}{proved=annex:p04-universal.2}{Adequacy of the numbering: enumeration theorem}
\begin{enumerate}
  \item[(a)] Every \(\varphi_e\) is partial computable.
  \item[(b)] For every \(k\ge1\) and every partial computable
        \(f:\N^k\rightharpoonup\N\), there is an index \(e\) such that
        \[
          \varphi_e(\vec x)\simeq f(\vec x)
          \qquad\text{for every }\vec x\in\N^k,
        \]
        under Convention~\ref{def:p04-program-numbering}'s arity-\(k\)
        application notation.
\end{enumerate}
\end{hbtheorem}

\begin{hbcorollary}{thm:p04-numbering-partial-recursive}{proved=annex:p04-universal.9}{Range of the local numbering}
For every \(k\geq1\) and partial map \(f:\N^k\rightharpoonup\N\), the
following are equivalent:
\begin{enumerate}
  \item some index \(e\) in Convention~\ref{def:p04-program-numbering}
        satisfies
        \[
          \varphi_e(\vec x)\simeq f(\vec x)
          \qquad\text{for every }\vec x\in\N^k;
        \]
  \item \(f\) is partial recursive in the closure sense of
        Definition~\ref{def:p04-partial-recursive}.
\end{enumerate}
\end{hbcorollary}

Theorem~\ref{thm:p04-enumeration} is the exact R-005 repair.  Its in-document
proof is primary because normalization, the binary-alphabet validity condition,
and the wrapper are bespoke.  Kleene's enumeration and normal-form treatment
for standard numberings is parallel reading \cite{Kleene1952}.

\begin{hbtheorem}{thm:p04-universal-computation}{sketched=annex:p04-universal.3}{Universal computation}
There is a partial computable function \(\mathsf{Univ}(e,x)\) such that
\[
  \mathsf{Univ}(e,x)\simeq\varphi_e(x)
\]
for all \(e,x\in\N\).
\end{hbtheorem}

Turing's original universal-machine construction is the historical source
\cite{Turing1937}.

\begin{hbtheorem}{thm:p04-parameter}{cited=Kleene1952}{Parameter theorem}
For every \(m,n\ge1\), there is a total computable function
\(\mathsf{par}_n^m:\N^{m+1}\to\N\) such that, for every index \(e\), every
\(\vec a\in\N^m\), and every \(\vec x\in\N^n\),
\[
  \varphi_{\mathsf{par}_n^m(e,\vec a)}(\vec x)
  \simeq\varphi_e(\vec a,\vec x).
\]
\end{hbtheorem}

The compiler encoded by \(\mathsf{par}_n^m\) decodes the runtime input
\(\vec x\), computes \(\operatorname{code}_{m+n}(\vec a,\vec x)\), writes its
binary word, and simulates \(e\).  The previously fixed tuple algorithms,
state renumbering, and finite-table concatenation make the index transformation
total and computable; \cite{Kleene1952} supplies the complete theorem for
effective numberings.

\section{Many-one reductions}
\label{sec:p04-many-one}

\begin{hbdefinition}{def:p04-many-one-reduction}{Many-one reduction}
For \(A,B\subseteq\N\), a \emph{many-one reduction} from \(A\) to \(B\) is a
total computable \(f:\N\to\N\) such that
\[
  x\in A\quad\Leftrightarrow\quad f(x)\in B
  \qquad\text{for every }x\in\N.
\]
Write \(A\leq_m B\) when one exists.
The sets are \emph{many-one equivalent}, written \(A\sim_m B\), exactly when
both \(A\leq_m B\) and \(B\leq_m A\); unlike one-way reducibility, this is
the relation of mutual many-one reducibility.
\index{many-one reduction}
\index{many-one equivalence}
\end{hbdefinition}

\begin{hblemma}{thm:p04-reduction-preservation}{proved=annex:p04-universal.4}{Preservation under many-one reduction}
If \(A\leq_m B\) and \(B\) is decidable, then \(A\) is decidable.  If
\(A\leq_m B\) and \(B\) is computably enumerable, then \(A\) is computably
enumerable.
\end{hblemma}

\begin{hbwarning}{def:p04-reduction-direction}{Direction of a reduction}
To prove that \(B\) is at least as hard as \(A\), construct \(A\leq_m B\).
The reverse reduction does not establish that lower bound.
\end{hbwarning}

\section{The self-application domain problem}
\label{sec:p04-self-application}

\begin{hbdefinition}{def:p04-self-application-domain}{Self-application domain set}
The \emph{self-application domain set} for the fixed numbering is
\[
  K=\{e\in\N:\varphi_e(e)\text{ is defined}\}.
\]
Here ``defined'' retains Convention~\ref{def:p04-program-numbering}'s exact
accepting-with-valid-output meaning.
\index{self-application domain set}
\end{hbdefinition}

\begin{hbtheorem}{thm:p04-halting-undecidable}{proved=annex:p04-universal.5}{Undecidability of the self-application domain set}
The set \(K\) is computably enumerable but not decidable.
\end{hbtheorem}

\begin{hbcorollary}{thm:p04-k-complement}{proved=annex:p04-universal.6}{The complementary domain set}
The complement \(\N\setminus K\) is not computably enumerable.
\end{hbcorollary}

\begin{hbdefinition}{def:p04-binary-application-domain}{Binary application domain set}
The \emph{binary application domain set} is
\[
  H=\{\operatorname{pair}_{\N}(e,x):e,x\in\N
       \text{ and }\varphi_e(x)\text{ is defined}\}.
\]
\index{binary application domain set}
\end{hbdefinition}

\begin{hbcorollary}{thm:p04-binary-application-equivalence}{proved=annex:p04-universal.10}{Binary and self-application domain problems}
The sets \(H\) and \(K\) are many-one equivalent:
\[
  H\sim_m K.
\]
In particular, \(H\) is undecidable.
\end{hbcorollary}

\section{The recursion theorem}
\label{sec:p04-recursion-theorem}

\begin{hbtheorem}{thm:p04-recursion-theorem}{sketched=annex:p04-universal.7}{Kleene recursion theorem}
If \(f:\N\to\N\) is total computable, then there is an index \(e\) such
that
\[
  \varphi_e\simeq\varphi_{f(e)}.
\]
\end{hbtheorem}

\begin{hbexample}{def:p04-self-reference-without-inspection}{Self-reference without a self-inspection primitive}
If a total computable transformation \(f\) adds a fixed output header to the
behavior of the program whose index it receives, the recursion theorem
supplies \(e\) whose behavior is that of \(f(e)\).  It asserts
\(\varphi_e\simeq\varphi_{f(e)}\), not \(e=f(e)\).
\end{hbexample}

\section{Rice's theorem}
\label{sec:p04-rice}

\begin{hbdefinition}{def:p04-extensional-index-property}{Extensional index property}
A set \(A\subseteq\N\) is \emph{extensional} iff, for all \(e,d\in\N\),
\[
  \varphi_e\simeq\varphi_d
  \quad\text{implies}\quad
  (e\in A\Leftrightarrow d\in A).
\]
It is \emph{nontrivial} iff it is neither empty nor all of \(\N\).
Thus extensionality classifies indices only through the partial functions they
compute, whereas nontriviality excludes the two constant classifications; both
distinctions are used in Theorem~\ref{thm:p04-rice}.
\index{extensional index property}
\index{nontrivial index property}
\end{hbdefinition}

\begin{hbtheorem}{thm:p04-rice}{proved=annex:p04-universal.8}{Rice's theorem}
Every nontrivial extensional set of indices of unary partial computable
functions is undecidable.
\end{hbtheorem}

\begin{hbwarning}{def:p04-rice-scope}{Semantic scope of Rice's theorem}
Rice's theorem concerns semantic properties of the partial function computed
by a program.  A syntactic property such as ``the source contains ten states''
can be decidable and is not extensional.
\end{hbwarning}

% SOURCE ACCOUNTING -- RETIRED EXERCISE ex:universal-reductions-01.
% The source exercise applying Rice's theorem to totality, range containing
% zero, and nonempty-domain constant partial functions is deliberately omitted
% because the handbook has no exercise layer.  The third property's nonempty
% domain proviso is recorded here so the retirement does not erase its scope.
